\section{Evaluation}
\label{sec:evaluation}
\subsection{Protocol and evidence classes}
We distinguish a released baseline execution, an in-library layout microbenchmark, and a cost model. The first supplies observed wall-clock and bootstrap events; the second isolates arithmetic layout counts under our evaluator; the third supplies uninstrumented baseline quantities. A counter-derived normalization is not a separate execution. Historical generated tables retain the working name \emph{hekit} for \sys{} and its crate names.

The headline uses MNIST reduced to $196$ inputs and the released full training setting. FULL-1 has completed at $B=64$ with one initial-weight seed; FULL-2 at the released baseline batch remains pending. FULL-1 is partial full-setting evidence, not a substitute for the FULL-2 accuracy gate. The completed faithful encrypted subset uses $784$ inputs, the baseline's sample indices and initial weights, and its training conventions. Output width is padded to $16$, including the unused classes in the released loss; replacing this with a width-$10$ loss is not assumed equivalent. Earlier ``legacy'' runs use a different learning rate, bias policy, and initialization. Their counters and timing remain useful, but their accuracy is not comparable to the released baseline. The last subset batch is partial, affecting depth, refresh scheduling, and the effective sum-versus-mean learning-rate conversion.

The paired server measurements use the same CPU and thread cap, with different CPU parallel runtimes. Jobs are not core-pinned, and some rows overlap unrelated shared-host load. Client work and setup scopes also differ: baseline step time includes batch encryption and prediction decryption, whereas our update time excludes those separately reported client operations; both exclude key setup. Thus even a same-host ratio is not a controlled architectural speedup. We do not pool CPU and GPU samples. Encrypted rows have a single seed unless the table states otherwise; within-run iterations and Criterion samples are not independent seeds.

\subsection{RQ1: key switches and wall-clock}
% generated by tools/paper_tables.py from results/; do not edit
\begin{table*}[t]
\centering
\scriptsize
\caption{RQ1: key switches per sample, four ways (ratio hekit / ReBoot at B = 48)}
\label{tab:rq1_ks}
\resizebox{\linewidth}{!}{%
\begin{tabular}{lrrrrrrrrrrrrr}
\toprule
Run & Input & Placement & B & Config & Samples & hekit KS/sample & Boots & ReBoot non-boot & (a) & (b) & (c) & (d) & Ways $\leq$ 0.70 \\
\midrule
\multicolumn{14}{l}{\textit{196 inputs, ReBoot full setting, 4800 samples x 10 epochs, seed 0 (headline)}} \\
FULL-1 & 196 & WN & 64 & faithful & 48000 & 18.272 & 1500 & 37 (model) & 0.232 & 0.387 & 0.395 & 0.400 & abcd \\
FULL-2 & 196 & WN & 48 in 64 & faithful & 48000 & 28.990 & 4000 & 37 (model) & 0.368 & 0.781 & 0.627 & 0.534 & acd \\
\midrule
\multicolumn{14}{l}{\textit{784 inputs, ReBoot encrypted-script subset, 1 epoch (supplementary)}} \\
F1 & 784 & WN & 32 & faithful & 1008 & 33.488 & 68 & 35.0625 (Lazy, measured) & 0.435 & 0.779 & 0.756 & 0.742 & a \\
F2 & 784 & WN & 64 & faithful & 1008 & 17.885 & 36 & 35.0625 (Lazy, measured) & 0.233 & 0.414 & 0.404 & 0.397 & abcd \\
R-41-3 & 784 & WN & 32 & legacy (R-M3-12) & 1024 & 32.531 & 64 & 35.0625 (Lazy, measured) & 0.423 & 0.741 & 0.734 & 0.730 & a \\
R-41-4 & 784 & WN & 64 & legacy (R-M3-12) & 1024 & 17.172 & 32 & 35.0625 (Lazy, measured) & 0.223 & 0.382 & 0.388 & 0.391 & abcd \\
R-41-1 & 784 & AU & 32 & legacy (R-M3-12) & 1024 & 39.469 & 128 & 35.0625 (Lazy, measured) & 0.513 & 1.149 & 0.891 & 0.730 & a \\
R-41-2 & 784 & AU & 64 & legacy (R-M3-12) & 1024 & 20.641 & 64 & 35.0625 (Lazy, measured) & 0.268 & 0.586 & 0.466 & 0.391 & abcd \\
R-41-5 & 784 & AA & 32 & legacy (R-M3-12) & 1024 & 36.000 & 96 & 35.0625 (Lazy, measured) & 0.468 & 0.945 & 0.812 & 0.730 & a \\
\midrule
\multicolumn{14}{l}{\textit{Interactive hekit, client refresh, no bootstraps (separate row, not way (d))}} \\
interactive & 784 & --- & 32 & legacy (M2) & 1024 & 19.781 & 0 & 35.0625 (Lazy, measured) & --- & --- & --- & --- & --- \\
\bottomrule
\end{tabular}}
\par\smallskip\begin{minipage}{\linewidth}\scriptsize
Independent interactive reference: ratio 0.564 against ReBoot Lazy; client refresh, not NI-derived way (d) and not a four-way pass result.\par
KS = rotations + relinearizations + SSE key switches (OpCounter). hekit bootstrap = 111 KS (measured, setup.boot\_ks). ReBoot bootstrap = 502 KS and 196-input non-boot = 37 KS/sample are M0 cost-model values (OpenFHE does not count key switches); ReBoot bootstraps 4 ciphertexts per update at B = 48 (measured on the official subset run: 84 bootstraps / 21 steps).\par
(a) own bootstrap costs; (b) both 502; (c) both 111; (d) bootstraps excluded on both sides. (b)-(d) are derived from the measured counters, not separate executions. Criterion: ratio $\leq$ 0.70 (research plan \S{}4.3.2).\par
784 rows are the supplementary subset rows (R-M3-13): ReBoot non-boot = Lazy 35.06, the in-library measurement of ReBoot's layout (Task 33, baseline c). Legacy rows use the pre-R-M3-12 config; their counts equal the faithful config but their accuracy is not comparable.\par
\end{minipage}
\end{table*}

% generated by tools/paper_tables.py from results/; do not edit
\begin{table*}[t]
\centering
\scriptsize
\caption{Wall-clock on the same CPU (quadruped-gpu, i9-13900KF, 12 threads)}
\label{tab:wallclock_same_cpu}
\resizebox{\linewidth}{!}{%
\begin{tabular}{lrrrrrrrrr}
\toprule
Run & Input & B & Samples & Train s & s/sample & vs ReBoot & Boot s/sample & Non-boot s/sample & Threads \\
\midrule
\multicolumn{10}{l}{\textit{196 inputs, full setting (headline)}} \\
FULL-1 & 196 & 64 & 48000 & pending (R-41-FULL-1) &  & ReBoot 196 full: not measured &  &  &  \\
FULL-2 & 196 & 48 in 64 & 48000 & pending (R-41-FULL-2) &  & ReBoot 196 full: not measured &  &  &  \\
\midrule
\multicolumn{10}{l}{\textit{784 inputs, ReBoot encrypted-script subset}} \\
ReBoot official (OpenFHE) & 784 & 48 & 1008 & 3599.9 & 3.571 & 1 & 2.585 & 0.986 & OpenMP 12 \\
F1 & 784 & 32 & 1008 & 3016.7 & 2.993 & 0.838 & 0.989 & 2.003 & rayon 12 \\
F2 & 784 & 64 & 1008 & 1607.9 & 1.595 & 0.447 & 0.522 & 1.073 & rayon 12 \\
R-41-3 & 784 & 32 & 1024 & 2773.6 & 2.709 & 0.758 & 0.871 & 1.838 & rayon 12 \\
R-41-4 & 784 & 64 & 1024 & 1434.8 & 1.401 & 0.392 & 0.435 & 0.967 & rayon 12 \\
R-41-1 & 784 & 32 & 1024 & 4339.1 & 4.237 & 1.187 & 1.736 & 2.501 & rayon 12 \\
R-41-2 & 784 & 64 & 1024 & 2228.9 & 2.177 & 0.609 & 0.867 & 1.310 & rayon 12 \\
R-41-5 & 784 & 32 & 1024 & 3606.6 & 3.522 & 0.986 & 1.316 & 2.206 & rayon 12 \\
\bottomrule
\end{tabular}}
\par\smallskip\begin{minipage}{\linewidth}\scriptsize
ReBoot time = sum of its own step\_time (encrypt batch, forward/backward, decrypt prediction, weight bootstrap; setup excluded). hekit time = time\_s.train\_total (all updates including bootstraps; key generation and bootstrapper setup excluded). Present rows fail closed unless the binary's provenance names the same host and thread count as the ReBoot run.\par
Shared machine: other users' jobs (SLAM, \textasciitilde{}4 cores) ran during some rows; see docs/m3-report.md \S{}3.2. ReBoot's full 196-input setting was not run (estimated 47.6 h), so the headline has no same-CPU wall-clock ratio.\par
\end{minipage}
\end{table*}

Let $K$ be measured total key switches, $H$ bootstrap events, and $S$ examples processed. Let $c_h$ be our measured per-bootstrap cost and $c_r$ the baseline bootstrap model; let $R$ denote the baseline non-bootstrap count per example, $r$ its refreshes per update, and $B_r$ its own batch. The four reported ratios are
\begin{align}
 \rho_a&=\frac{K/S}{R+rc_r/B_r},&
 \rho_b&=\frac{(K-c_hH+c_rH)/S}{R+rc_r/B_r},\nonumber\\
 \rho_c&=\frac{K/S}{R+rc_h/B_r},&
 \rho_d&=\frac{(K-c_hH)/S}{R}.
 \label{eq:ratios}
\end{align}
Convention (a) uses each side's own bootstrap cost. Convention (b) charges both the baseline cost; (c) charges both our cost; (d) excludes bootstrap arithmetic on both sides. The interactive comparison is separate from (d): omitting bootstrap arithmetic does not remove non-interactive compact packing and unpacking. ReBoot remains at its own batch in the denominator; a hypothetical same-batch denominator would answer a different question.

The baseline's bootstrap cost remains a model because the measured OpenFHE execution does not expose key-switch counters. Its observed bootstrap frequency is used, not replaced with our schedule. The subset non-bootstrap reference is the in-library Lazy layout measurement; the full-setting reference is still a model. These boundaries are visible in Table~\ref{tab:rq1_ks}.

% Generated by tools/figs.py; do not edit. See rq1_ratios.sources.json.
\begin{figure}[t]
\centering
\setlength{\unitlength}{1pt}
\begin{picture}(240,143)
\put(28,32){\line(1,0){204}}
\put(28,32){\line(0,1){82}}
\put(5,124){\makebox(225,10){\footnotesize KS ratio to ReBoot (lower is better)}}
\put(25,32.00){\line(1,0){3}}
\put(0,29.00){\makebox(22,7)[r]{\tiny 0}}
\put(25,52.00){\line(1,0){3}}
\put(0,49.00){\makebox(22,7)[r]{\tiny 0.25}}
\put(25,72.00){\line(1,0){3}}
\put(0,69.00){\makebox(22,7)[r]{\tiny 0.5}}
\put(25,92.00){\line(1,0){3}}
\put(0,89.00){\makebox(22,7)[r]{\tiny 0.75}}
\put(25,112.00){\line(1,0){3}}
\put(0,109.00){\makebox(22,7)[r]{\tiny 1}}
\multiput(28,88.00)(5,0){41}{\line(1,0){2}}
\put(176,91.00){\makebox(55,7)[r]{\tiny threshold 0.7}}
\put(42,32){\color[HTML]{315C8C}\rule{14pt}{34.800pt}}
\put(38,68.800){\makebox(22,7){\tiny 0.435}}
\put(42,21){\makebox(14,7){\scriptsize (a)}}
\put(61,32){\color[HTML]{487F9C}\rule{14pt}{62.320pt}}
\put(57,96.320){\makebox(22,7){\tiny 0.779}}
\put(61,21){\makebox(14,7){\scriptsize (b)}}
\put(80,32){\color[HTML]{6C9B8B}\rule{14pt}{60.480pt}}
\put(76,94.480){\makebox(22,7){\tiny 0.756}}
\put(80,21){\makebox(14,7){\scriptsize (c)}}
\put(99,32){\color[HTML]{A2A76B}\rule{14pt}{59.360pt}}
\put(95,93.360){\makebox(22,7){\tiny 0.742}}
\put(99,21){\makebox(14,7){\scriptsize (d)}}
\put(35,6){\makebox(90,10){\footnotesize F1: B = 32}}
\put(145,32){\color[HTML]{315C8C}\rule{14pt}{18.640pt}}
\put(141,52.640){\makebox(22,7){\tiny 0.233}}
\put(145,21){\makebox(14,7){\scriptsize (a)}}
\put(164,32){\color[HTML]{487F9C}\rule{14pt}{33.120pt}}
\put(160,67.120){\makebox(22,7){\tiny 0.414}}
\put(164,21){\makebox(14,7){\scriptsize (b)}}
\put(183,32){\color[HTML]{6C9B8B}\rule{14pt}{32.320pt}}
\put(179,66.320){\makebox(22,7){\tiny 0.404}}
\put(183,21){\makebox(14,7){\scriptsize (c)}}
\put(202,32){\color[HTML]{A2A76B}\rule{14pt}{31.760pt}}
\put(198,65.760){\makebox(22,7){\tiny 0.397}}
\put(202,21){\makebox(14,7){\scriptsize (d)}}
\put(138,6){\makebox(90,10){\footnotesize F2: B = 64}}
\end{picture}
\caption{Supplementary subset RQ1 ratios (one seed per row), generated from the same cells as Table~\ref{tab:rq1_ks}. Bars (a)--(d) denote the four accounting conventions in Section~\ref{sec:evaluation}. These are not equal-security or full-setting results.}
\Description{Grouped bars for the two faithful subset batches under four key-switch accounting conventions. The larger batch is below the threshold in all four; the smaller batch only under own bootstrap costs. One seed per row; not an equal-security comparison.}
\label{fig:rq1-ratios}
\end{figure}

The faithful larger-batch subset row meets the registered count threshold under every reported convention, while the smaller batch meets it only under the own-bootstrap convention. The generated figure makes that dependence visible. The corresponding larger-batch timing is also favorable on the shared CPU (Table~\ref{tab:wallclock_same_cpu}), but its subset accuracy drops relative to the baseline, including in the plaintext twin. Consequently this observation is a tradeoff, not a successful accuracy-matched RQ1 result. FULL-1 also meets the count threshold under all four conventions. Its completed training time has no measured full-setting ReBoot CPU denominator. The pending FULL-2 accuracy gate, missing baseline timing, and security mismatch below prevent a final verdict.

\subsection{RQ2: placement and depth}
The AU, WN, and AA subset rows in Tables~\ref{tab:rq1_ks} and~\ref{tab:wallclock_same_cpu} provide the placement comparison. For matched legacy eMLP-1 settings, WN reduces bootstrap events and total time relative to AU. AA leaves the same lazy weight-refresh frequency and adds activation refreshes, so it does not improve this shallow workload. Its benefit is a feasible depth budget for some deeper polynomial circuits, checked separately as a structural test rather than an observed training-accuracy gain.

The corrected structural predictor equals the measured counters at every completed update. This validates accounting for the executed paths; it does not validate the baseline bootstrap model or a convergence claim. The initial cost model missed replication, key-reuse, bias, and schedule-related terms, and its optimistic interactive estimate is superseded by that explicit accounting. FULL-1 extends exact schedule agreement to its complete registered update sequence and records accumulated weight and bias error against its same-batch plaintext twin. That observed error is larger than in the earlier subset checks; close final accuracy does not imply a uniform small parameter error. FULL-2 and guarantees across seeds or arbitrary trajectories remain unmeasured.

\subsection{Accuracy, local blocks, and security}
% generated by tools/paper_tables.py from results/; do not edit
\begin{table*}[t]
\centering
\scriptsize
\caption{Test accuracy against ReBoot (R-M3-15)}
\label{tab:accuracy_vs_reboot}
\resizebox{\linewidth}{!}{%
\begin{tabular}{lrrrrrrr}
\toprule
Row & Setting & Seeds & Test acc (\%) & Twin (\%) & Reference (\%) & $\Delta$ (pt) & Gate \\
\midrule
\multicolumn{8}{l}{\textit{Full setting: 196 inputs, 4800 samples x 10 epochs, B = 48 (the $\pm$0.5 pt gate)}} \\
ReBoot released code, plaintext (recomputed twin) & full & 10 & 90.65 $\pm$ 0.27 (n = 10) & --- & --- & --- & reference \\
hekit plaintext twin (ReBoot weights, order, lr) & full & 10 & 90.65 $\pm$ 0.27 (n = 10) & --- & 90.65 $\pm$ 0.27 (n = 10) & +0.00 & twin check \\
hekit encrypted, B = 48 in 64 lanes, WN & full & 1 (seed 0) & pending (R-41-FULL-2) & pending (R-41-FULL-2) & 90.41 &  & pending \\
hekit encrypted, B = 64, WN & full & 1 (seed 0) & 90.36 & 90.37 & 90.41 (seed 0) & -0.05 & not the gate (B $\neq$ 48) \\
\midrule
\multicolumn{8}{l}{\textit{eMLP-2 / eMLP-3 block MLPs, full setting, seed 0}} \\
hekit plaintext twin emlp2 (blocks 64) & full & 1 (seed 0) & 91.56 & --- & 91.56 (ReBoot twin) & +0.00 & twin check \\
hekit plaintext twin emlp3 (blocks 128,64) & full & 1 (seed 0) & 91.75 & --- & 91.75 (ReBoot twin) & +0.00 & twin check \\
\midrule
\multicolumn{8}{l}{\textit{784 inputs, ReBoot encrypted-script subset, 1008 samples, 1 epoch (not a gate, R-M3-15)}} \\
ReBoot official encrypted (OpenFHE) & subset, B = 48 & 1 (seed 0) & 75.06 & --- & --- & --- & reference \\
hekit encrypted WN, B = 32 & subset & 1 (seed 0) & 75.77 & 75.77 & 75.06 & +0.71 & not a gate \\
hekit encrypted WN, B = 64 & subset & 1 (seed 0) & 73.58 & 73.59 & 75.06 & -1.48 & not a gate \\
hekit encrypted eMLP-2, B = 48 (196) & subset 1008 & 1 & pending (R-52-1) & pending (R-52-1) &  &  & twin check \\
hekit encrypted eMLP-3, B = 48 (196) & subset 1008 & 1 & pending (R-52-2) & pending (R-52-2) &  &  & twin check \\
\midrule
\multicolumn{8}{l}{\textit{Reference only: the paper value and the released code at the paper setting (R-REB-3)}} \\
ReBoot paper, Table 4 (plaintext) & 60k samples & 10 (paper) & 94.61 $\pm$ 0.27 & --- & --- & --- & not reproducible from the released defaults \\
ReBoot released code, paper setting, constant lr, epoch 50 & 60k samples, 50 epochs & 10 & 95.15 $\pm$ 0.19 (n = 10) & --- & 94.61 & +0.54 & reference \\
ReBoot released code, paper setting, constant lr, best epoch (test-selected) & 60k samples, 50 epochs & 10 & 95.35 $\pm$ 0.10 (n = 10) & --- & 94.61 & +0.74 & reference \\
ReBoot released code, paper setting, ReBoot's scheduler, epoch 50 & 60k samples, 50 epochs & 10 & 95.25 $\pm$ 0.08 (n = 10) & --- & 94.61 & +0.64 & reference \\
ReBoot released code, paper setting, ReBoot's scheduler, best epoch (test-selected) & 60k samples, 50 epochs & 10 & 95.46 $\pm$ 0.09 (n = 10) & --- & 94.61 & +0.85 & reference \\
\bottomrule
\end{tabular}}
\par\smallskip\begin{minipage}{\linewidth}\scriptsize
FULL-1 final parameter error: 0.00682657 (maximum absolute difference across weights and biases against its same-batch plaintext twin at training end). This is a trajectory diagnostic, not a per-bootstrap precision or accuracy gate.\par
R-M3-15: the $\pm$0.5 pt gate compares hekit encrypted B = 48 (FULL-2) with ReBoot's released code at its full setting, seed 0. Subset rows are 1 epoch and are not a gate. The paper value and R-REB-3 are reported with this note only.\par
\end{minipage}
\end{table*}

% generated by tools/paper_tables.py from results/; do not edit
\begin{table*}[t]
\centering
\scriptsize
\caption{Plaintext dataset reproduction against ReBoot Table 4 (Task 51)}
\label{tab:plain_datasets}
\resizebox{\linewidth}{!}{%
\begin{tabular}{lrrrrrrrr}
\toprule
Dataset & Config & Epochs & Seeds & Fixed final epoch (\%) & Best epoch, test-selected (\%) & ReBoot Table 4 (\%) & $\Delta$ fixed (pt) & Within 1 pt \\
\midrule
MNIST 14x14 & paper setting & 50 & 10 & 95.00 $\pm$ 0.18 (n = 10) & 95.33 $\pm$ 0.15 (n = 10) & 94.61 $\pm$ 0.27 & +0.39 & yes \\
Fashion-MNIST 14x14 & paper setting & 50 & 10 & 85.22 $\pm$ 0.17 (n = 10) & 85.53 $\pm$ 0.12 (n = 10) & 85.12 $\pm$ 0.20 & +0.10 & yes \\
Breast Cancer & fixed split, min-max (primary) & 100 & 10 & 96.40 $\pm$ 1.27 (n = 10) & 99.30 $\pm$ 0.37 (n = 10) & 99.03 $\pm$ 1.78 & -2.63 & no \\
Breast Cancer & per-run split (sensitivity) & 100 & 10 & 95.26 $\pm$ 1.61 (n = 10) & 97.89 $\pm$ 0.94 (n = 10) & 99.03 $\pm$ 1.78 & -3.77 & no \\
Breast Cancer & z-score (diverges) & 100 & 10 & diverged (10/10 seeds) & --- & 99.03 $\pm$ 1.78 & --- & no \\
\bottomrule
\end{tabular}}
\par\smallskip\begin{minipage}{\linewidth}\scriptsize
Fixed final epoch is the primary figure. The best epoch is selected on the test set and is optimistic; it is shown only because ReBoot does not state its epoch count. Seeds: the plain binary's own seeds (not ReBoot's initial weights).\par
Fashion-MNIST and Breast Cancer are plaintext twins only (no encrypted runs).\par
\end{minipage}
\end{table*}

% generated by tools/paper_tables.py from results/; do not edit
\begin{table*}[t]
\centering
\scriptsize
\caption{Security estimates (lattice-estimator; log2 rop, minimum over attacks)}
\label{tab:security}
\resizebox{\linewidth}{!}{%
\begin{tabular}{lrrrrrr}
\toprule
Parameters & N & log PQ & Secret & BDGL16 & MATZOV & >= 128 (both) \\
\midrule
hekit (a): n16\_boot main key & 2\textasciicircum{}16 & 1732 (1731.998) & dense ternary (P($\pm$1) = 1/4 each, h $\approx$ N/2) & 131.5 & 129.3 & yes \\
hekit (b): SSE key s -> s' (h' = 32) & 2\textasciicircum{}16 & 294 (q\_0$\cdot$P) & sparse ternary, h' = 32 & 133.8 & 134.2 & yes \\
ReBoot main key (R-REB-4: OpenFHE's own chain) & 2\textasciicircum{}16 & 1730 (1729.585) & sparse ternary, h = 192 (97 / 95) & 121.8 & 121.5 & no \\
\bottomrule
\end{tabular}}
\par\smallskip\begin{minipage}{\linewidth}\scriptsize
ReBoot row: the MITM tier (operation counts; about 2\textasciicircum{}104-2\textasciicircum{}110 memory). Without MITM attacks ReBoot is at 127.8 / 126.2 (BDGL16 / MATZOV), also below 128.\par
Source: the summary table of docs/security.md (the estimator outputs are recorded in that doc; no JSON exists). Same lattice-estimator models for every row; m = N; Core-SVP is in the doc for information only. Equal N does not mean equal security: no equal-security claim.\par
\end{minipage}
\end{table*}

Table~\ref{tab:accuracy_vs_reboot} separates the released full-setting reference, subset results, and published paper accuracy. The released defaults and paper-setting reconstruction are different protocols. We do not choose whichever reference makes the result look best. The registered tolerance is tested by FULL-2 against the captured released full-setting seed-zero reference; that encrypted row has not completed. FULL-1 reports its larger-batch accuracy and same-batch twin separately, without substituting either for this gate. Matching a plaintext twin is necessary to diagnose arithmetic fidelity, but is insufficient to establish baseline-equivalent generalization.

Plaintext dataset experiments (Table~\ref{tab:plain_datasets}) provide broader optimization checks. MNIST and Fashion-MNIST reproduce the stated comparison tolerance at the fixed final epoch. Breast Cancer does not: the fixed-epoch min-max result misses its target, and selecting the best test epoch is optimistic leakage from model selection. The z-score setting diverges for the polynomial activation. Fashion-MNIST and Breast Cancer results are plaintext twins, not encrypted training measurements. The local-block plaintext twins reproduce the released eMLP-2/eMLP-3 values for the recorded seed; their encrypted subset rows are pending. The block experiments do not extend the current headline claim to deeper networks.

Table~\ref{tab:security} applies common estimator models to the actual exposed key moduli and secret distributions. The evaluated main and SSE instances of our preset exceed the chosen operation-count target under those models; the released ReBoot preset does not. This is an estimate over the documented attack search, not a security proof. The smallest sparse-secret values use attacks with extreme memory requirements, and even non-MITM minima have substantial or unreported memory demands. These are not practical attack demonstrations. Equal $N$ and a requested library security label do not establish equal security. Results must remain conditional until either a security-different comparison is explicitly accepted or an independently estimated and remeasured alternative baseline is available.

\subsection{RQ3: device operations and ablation}
% generated by tools/paper_tables.py from results/; do not edit
\begin{table*}[t]
\centering
\scriptsize
\caption{GPU vs CPU vs HEonGPU, n16\_train (R-47-2, criterion medians, ms)}
\label{tab:gpu_vs_cpu_ops}
\resizebox{\linewidth}{!}{%
\begin{tabular}{lrrrrrrr}
\toprule
Operation & Level & hekit GPU & CPU 1 thr & CPU 12 thr & GPU speed-up (1 / 12 thr) & HEonGPU (RTX 4090) & GPU / HEonGPU \\
\midrule
add & 31 & 0.111 & 15.674 & 14.611 & 141x / 132x & none published & --- \\
relinearize & 31 & 4.671 & 639.639 & 240.070 & 137x / 51x & 2.582 (l = 36, I = 4 primes per digit) & 1.81x (nominal) \\
mul + relin + rescale & 31 & 5.823 & 725.209 & 337.731 & 125x / 58x & none published & --- \\
rotate(1) & 31 & 3.990 & 622.046 & 223.880 & 156x / 56x & none published & --- \\
hoisted rotate x16 & 31 & skipped: needs 5580 MiB (2377 MiB free) & 3765.068 & 689.471 & --- & none published & --- \\
hoisted rotate x16 (capped keys) & 9 & 4.984 & 1005.331 & 172.870 & 202x / 35x & none published & --- \\
matvec\_enc d = 64 & 31 & skipped: needs 7276 MiB (1929 MiB free) & 12938.440 & 2325.173 & --- & none published & --- \\
matvec\_enc d = 64 (capped keys) & 9 & 16.120 & 2523.223 & 547.762 & 157x / 34x & none published & --- \\
matvec\_enc d = 64 (capped keys) & 5 & 9.509 & 1391.064 & 314.649 & 146x / 33x & none published & --- \\
matvec\_enc x4, 1 stream (sequential) & 5 & 34.994 & --- & --- & --- & none published & --- \\
matvec\_enc x4, 4 streams (untuned R-47-2) & 5 & 41.673 & --- & --- & --- & none published & --- \\
matvec\_enc x4, re-measured after the fix round & 5 & pending (R-47-3) & --- & --- & --- & --- & --- \\
\bottomrule
\end{tabular}}
\par\smallskip\begin{minipage}{\linewidth}\scriptsize
HEonGPU relinearize: Ozcan and Savas, GPU Implementations of Three Different Key-Switching Methods for Homomorphic Encryption Schemes, IACR ePrint 2025/124, Table 3: CKKS relinearization on a GeForce RTX 4090, log N = 16, 43-bit precision, Method II (hybrid), (I, l\textasciitilde{}, l) = (4, 40, 36), level 35. The previously cited Nov 2024 manuscript version date is unverified; ePrint landing metadata alone does not establish that manuscript date. Caveat: not the same parameters as n16\_train (36 vs 32 Q primes, 4 special primes; I = 4 primes per digit corresponds to alpha = 4, not four digits; n16\_train uses dnum 8); the ratio is nominal. Scaling key switching linearly in l gives about 2.30 ms at l = 32 and a ratio of about 2.0x (docs/m4-report.md section 1.1). ePrint 2024/1543 could not be read (HTTP 403), so other HEonGPU operation values are unconfirmed.\par
Capped keys: the input is dropped to level 9 or 5 and the keys are cut to that level (faithful for that level, not the top-level cost); the CPU uses full keys at the same level. The shared RTX 4090 had 2-3 GiB free, so full-key GPU rows that need more were skipped.\par
\end{minipage}
\end{table*}

% generated by tools/paper_tables.py from results/; do not edit
\begin{table*}[t]
\centering
\scriptsize
\caption{RQ3 ablation: one interactive eMLP-1 iteration under each EvalConfig}
\label{tab:rq3_ablation}
\resizebox{\linewidth}{!}{%
\begin{tabular}{lrrrrrrrrr}
\toprule
Source & Backend / level & hoist bsgs lazy & Rotations & Relins & KS/iter & Closed form & Step median (s) & Step / default & Max |$\Delta$W| vs default \\
\midrule
CPU, n16\_train, level 31 &  &  &  &  &  &  & pending (R-48-4) &  &  \\
GPU, n16\_train, level 7 &  &  &  &  &  &  & pending (R-48-2) &  &  \\
CPU, n16\_train, level 7 &  &  &  &  &  &  & pending (R-48-3) &  &  \\
CPU smoke (insecure N = 2\textasciicircum{}12: counts only) & cpu, level max & 1 1 1 & 171 & 62 & 233 & match & n/a (smoke) & n/a (smoke) & 0 \\
CPU smoke (insecure N = 2\textasciicircum{}12: counts only) & cpu, level max & 1 1 0 & 171 & 114 & 285 & match & n/a (smoke) & n/a (smoke) & 7.83e-09 \\
CPU smoke (insecure N = 2\textasciicircum{}12: counts only) & cpu, level max & 1 0 1 & 210 & 101 & 311 & match & n/a (smoke) & n/a (smoke) & 6.90e-09 \\
CPU smoke (insecure N = 2\textasciicircum{}12: counts only) & cpu, level max & 1 0 0 & 210 & 114 & 324 & match & n/a (smoke) & n/a (smoke) & 7.39e-09 \\
CPU smoke (insecure N = 2\textasciicircum{}12: counts only) & cpu, level max & 0 1 1 & 171 & 62 & 233 & match & n/a (smoke) & n/a (smoke) & 0 \\
CPU smoke (insecure N = 2\textasciicircum{}12: counts only) & cpu, level max & 0 1 0 & 171 & 114 & 285 & match & n/a (smoke) & n/a (smoke) & 7.83e-09 \\
CPU smoke (insecure N = 2\textasciicircum{}12: counts only) & cpu, level max & 0 0 1 & 210 & 101 & 311 & match & n/a (smoke) & n/a (smoke) & 6.90e-09 \\
CPU smoke (insecure N = 2\textasciicircum{}12: counts only) & cpu, level max & 0 0 0 & 210 & 114 & 324 & match & n/a (smoke) & n/a (smoke) & 7.39e-09 \\
\bottomrule
\end{tabular}}
\par\smallskip\begin{minipage}{\linewidth}\scriptsize
1. Definition of bsgs = false. The giant rotation is not shared; the weights keep the same pre-rotated (BSGS) layout and the same Galois keys. Every product with a nonzero giant shift is finished (relinearized) and rotated on its own; the products of the zero-shift block are summed as with BSGS. Per matvec of dimension d (offset 0): d - 1 rotations, of which only the baby - 1 rotations of x can be hoisted; with lazy\_relin 1 + d - baby relins, without it d (EvalConfig::matvec\_rotations, EvalConfig::matvec\_enc\_relins). At d = 16 (baby = giant = 4), as (hoistable rotations / individual rotations / relins): | form | hoistable rot. | individual rot. | relins (lazy) | relins (not lazy) | |---|---|---|---|---| | default BSGS | 3 | 3 | 4 | 16 | | bsgs = false (this library) | 3 | 12 | 13 | 16 | | textbook diagonal method (not implemented) | 15 | 0 | 1 | 16 | The textbook (Halevi-Shoup) diagonal method rotates x by all d - 1 steps (all hoistable) and sums the products lazily; it needs un-pre-rotated diagonals, i.e. another weight layout and gradient form, so it is not implemented. The all-off configuration has the same counts in both forms (15 rotations, 16 relins at d = 16).\par
2. Bias. Against the strongest naive baseline (the textbook method with hoisting and lazy relinearization), the cost of the bsgs = false rows overstates what BSGS gains: the measured BSGS gain is an upper bound. Conversely, the marginal gains of hoist and lazy\_relin measured inside bsgs = false configurations are lower bounds (only baby - 1 rotations can be hoisted and only the zero-shift block summed lazily: e.g. lazy\_relin saves 3 relins at d = 16 instead of 15).\par
3. hoist changes no OpCounter value. The counter does not distinguish hoisted from individual rotations, so the effect of hoist shows only in wall-clock time, never in key-switch counts.\par
4. bsgs does not touch the gradient's outer product (Dense::weight\_grad\_scaled, Dense::weight\_grad\_reduced): its (baby - 1) + (giant - 1) input rotations are what the pre-rotated weight layout requires, not a BSGS choice; only hoist changes how they run. The interactive gradient multiplies each diagonal with one Evaluator::mul (nothing to sum lazily), so lazy\_relin does not change it either.\par
The smoke rows (insecure N = 2\textasciicircum{}12, high Mac load) check counts and closed forms only; their times are not comparative.\par
\end{minipage}
\end{table*}

Table~\ref{tab:gpu_vs_cpu_ops} contains synchronized operation latencies, including capped-level rows and memory skips. The HEonGPU relinearization reference~\cite{gpuKeys2025} is nominal: its modulus count differs. A linear-in-limb adjustment is a cost-model extrapolation, not an equal-parameter measurement. The published table value is inherited from the archived operator report; the bibliographic metadata was checked for this draft, but a fresh PDF fetch was blocked. The later pool/submission fix has not been retimed.

Table~\ref{tab:rq3_ablation} deliberately leaves the server rows pending. The insecure small-parameter CPU smoke validates counts and checks only; its high-load timings are inadmissible for performance comparisons. The flag named \texttt{bsgs=false} keeps pre-rotated weights rather than implementing the strongest textbook diagonal baseline. Thus its BSGS gain would be an upper bound against that alternative, while the marginal hoisting and lazy-relinearization gains within those rows would be lower bounds. Hoisting changes no operation counter; the gradient outer product retains layout-required rotations independent of the BSGS flag. All four caveats travel with the table.

CPU and GPU training comparisons require the same optimizer, level, layout, and batch. Registered SGD fallback rows are compared only with the corresponding CPU SGD row. A resource skip is unmeasured; a device OOM is a failure. Test-process success cannot stand in for an actually executed production-size parity gate. GPU bootstrap precision, latency, allocation behavior, and peak memory are still unmeasured. HEonGPU's separate bootstrap reference~\cite{heongpuBench} has a different remaining-level budget and main-secret distribution; the fixed plan target is not a normalized speed-ratio result. RQ3's epoch-time question is therefore unanswered.
